\documentclass[12pt,a4paper]{report}

\usepackage[utf8]{inputenc}
\usepackage[T1]{fontenc}
\usepackage{geometry}
\usepackage{longtable}
\usepackage{array}
\usepackage{booktabs}
\usepackage{ragged2e}
\usepackage{enumitem}
\usepackage{hyperref}

\geometry{
    left=1.25in,
    right=1in,
    top=1in,
    bottom=1in
}

\begin{document}

\chapter{Design of the System}

This chapter describes the requirements, development approach, system
models, resource requirements, and architecture of the Agentic AI-Based
IT Procurement Intelligence Platform. It presents the functional and
non-functional requirements, the principal data flows, the system
layers, and the user-interface design. Diagram positions are identified
where the source report refers to figures that must be inserted or
regenerated.

\section{Requirement Engineering}

Requirement engineering involves identifying the capabilities expected
from the platform, documenting them clearly, and checking that they
remain consistent throughout development. In this project, it was
particularly important because the system is divided into services that
must interoperate through shared event formats, database structures,
REST response conventions, and authentication rules.

\subsection{Requirement Elicitation}

Requirements were identified from the following sources:

\begin{itemize}
    \item \textbf{Review of existing practices and literature:}
    The review in Chapter 2 and the problem analysis in Chapter 3
    identified recurring procurement challenges, including repeated
    manual data entry, unverified vendor details, duplicate invoices,
    delayed approvals, underused software licences, missed contract
    renewals, and infrequent vendor-risk reviews.

    \item \textbf{Stakeholder analysis:}
    The principal users are requesters, approvers and finance users,
    procurement/legal personnel, and system administrators. In the
    implemented demonstration application, the four roles are
    requester, approver, finance, and admin.

    \item \textbf{Guide and team discussions:}
    Discussions were used to prioritise features and distinguish the
    implemented scope from features dependent on external accounts or
    providers. Live GSTIN lookups, sanctions screening, and integration
    with a real e-signature provider may be affected by external
    dependencies or future scope.

    \item \textbf{Review of tools and datasets:}
    Document-processing tools, workflow engines, and relevant datasets
    were reviewed to understand their expected inputs, outputs, and
    practical limitations.
\end{itemize}

The gathered requirements were then grouped into functional and
non-functional requirements.

\subsection{Software Lifecycle Model}

The project followed an iterative and incremental development approach.
This approach was suitable because thresholds, approval rules, model
behaviour, and integration details required testing and refinement
during implementation.

The main practices were:

\begin{itemize}
    \item \textbf{Contract-first design:}
    Shared interfaces were agreed upon before parallel service
    development. These included the Kafka event envelope and topic
    payloads, the base database schema, the REST response envelope,
    and authentication middleware.

    \item \textbf{Service-wise increments:}
    Work was divided across document and vendor, approval and inventory,
    contract and risk, and platform/notification responsibilities.

    \item \textbf{Integration iterations:}
    The services were connected behind a common API gateway.
    Integration testing helped identify issues involving gateway
    routing, Kafka connectivity, and database migrations.

    \item \textbf{Continuous testing:}
    Service-level tests and end-to-end checks were used during
    development rather than postponing all testing until the end.

    \item \textbf{Repeatable deployment:}
    The project uses a \texttt{run.sh} script for bringing up the
    containerised stack, supporting consistent demonstrations and
    verification.
\end{itemize}

\textbf{Figure 4.1:} Iterative and incremental development cycle.
Insert or regenerate the diagram showing contract-first design,
service-wise increments, integration iteration, testing, and
demonstration, with guide feedback informing subsequent iterations.

\subsection{Requirement Analysis}

The requirements are grouped into functional requirements, which
describe what the system should do, and non-functional requirements,
which describe expected qualities and constraints.

\subsubsection{Functional Requirements}

\begin{longtable}{|p{0.8in}|p{3.8in}|p{1.5in}|}
\hline
\textbf{ID} & \textbf{Requirement} & \textbf{Responsible Service} \\
\hline
FR1 & Users can log in and receive access and refresh tokens. Protected
endpoints verify the token and user role, except specified health and
metrics endpoints. & Platform/Authentication \\
\hline
FR2 & Uploaded purchase orders, invoices, and quotations are scanned
for malware before storage. An upload is rejected if the scanner is
unavailable. & Document and Vendor \\
\hline
FR3 & Documents are classified as purchase orders, invoices, or
quotations. Relevant vendor, line-item, total, date, and document-number
fields are extracted. & Document and Vendor \\
\hline
FR4 & Extracted fields receive confidence scores. Results below the
configured threshold, with a default of 0.8, are sent for human review.
& Document and Vendor \\
\hline
FR5 & Vendor records are matched and duplicate documents are detected
using vendor, amount tolerance, and date-window checks.
& Document and Vendor \\
\hline
FR6 & Changes to vendor bank or payment details remain pending until
verified by a different user. & Document and Vendor \\
\hline
FR7 & Purchase requests follow approval routes based on spending tiers,
with service-level-agreement (SLA) timers and escalation.
& Approval and Inventory \\
\hline
FR8 & Approval and rejection decisions are applied through a durable
workflow, and the requester can see the resulting status.
& Approval and Inventory \\
\hline
FR9 & Inventory is reserved using locks to prevent concurrent requests
from reserving the same unit. Oversubscribed requests can be split into
immediate and back-ordered portions. & Approval and Inventory \\
\hline
FR10 & Licence utilisation is calculated, and a reclaim request can be
created when utilisation falls below a configured threshold.
& Approval and Inventory \\
\hline
FR11 & Contracts are generated from templates. Renewal type, notice
period, and end date are extracted, and reminders are scheduled for
60, 30, and 15 days before the notice deadline. & Contract and Risk \\
\hline
FR12 & The system sends contracts for signature and accepts signed-document
callbacks only when the HMAC signature is valid and the callback is not
a replay. & Contract and Risk \\
\hline
FR13 & Vendor risk is represented by a score and band, with contributing
factors listed and urgent alerts raised for a High band.
& Contract and Risk \\
\hline
FR14 & The risk model is monitored for drift; automatic retraining is
not performed. & Contract and Risk \\
\hline
FR15 & Vendor offboarding revokes access, flags active contracts for
review, and sets a data-retention flag rather than deleting records.
& Contract and Risk \\
\hline
FR16 & The notification service consumes events, sends templated email
alerts, prioritises urgent items, supports periodic digests for routine
items, and maintains a searchable notification log. & Notification \\
\hline
FR17 & The frontend provides role-specific navigation and dashboards,
purchase-request submission, approval inbox, document review, vendor,
contract and inventory views, and notification history. & Frontend \\
\hline
\end{longtable}

\subsubsection{Non-Functional Requirements}

\begin{longtable}{|p{1.5in}|p{4.8in}|}
\hline
\textbf{Category} & \textbf{Requirement} \\
\hline
Security & JWT-based authentication and backend role-based access
control; fail-closed malware scanning; HMAC-verified webhooks; and
dual control for payment-detail changes. \\
\hline
Reliability & Approval and renewal workflows are durable. Kafka
consumers are designed for idempotent processing, and webhook replays
are rejected. \\
\hline
Consistency & Endpoints follow a common response envelope, and shared
database changes use idempotent service migrations. \\
\hline
Performance & Uploads return promptly while document extraction and other
heavy processing run asynchronously. \\
\hline
Traceability & A correlation ID follows events and logs. Prometheus
metrics and Grafana dashboards support monitoring of request rate,
latency, errors, and service availability. \\
\hline
Portability & The stack is containerised and intended to run on macOS,
Linux, and Windows through WSL2. \\
\hline
Maintainability & The platform separates responsibilities by service,
uses shared contracts, and includes service-level and end-to-end tests. \\
\hline
Explainability & Risk scores display contributing factors, and
low-confidence extraction results are made available for human review. \\
\hline
\end{longtable}

\textbf{Constraints and assumptions:} Payment execution, ERP
integration, tax identifiers other than GSTIN and IFSC, and automatic
model retraining are outside the stated scope. Parts of the training
data are synthetic, and the availability of live external lookups
depends on free-tier quotas and external services.

\subsection{UML Diagrams and Data Flow Diagrams}

The diagrams in this section should be numbered consistently and
referenced in the text. Insert the actual diagrams as images or
recreate them in the final document.

\subsubsection{Use Case Diagram}

The use case diagram is presented in Figure 3.4 in Chapter 3. It
identifies the System Admin, Requester, Approver, and Procurement/Legal
Officer and groups the main activities into authentication and business
rules, purchase and inventory, vendor intelligence, contract and risk,
and notifications.

\subsubsection{Data Flow Diagram -- Level 0}

The Level 0 context diagram should represent the platform as one
process and show its interactions with the external entities identified
in the project. Insert the Level 0 diagram here and ensure its external
entities and labelled data flows are legible.

\subsubsection{Data Flow Diagram -- Level 1}

The Level 1 diagram decomposes the platform into the processes listed
in Table 4.2. Processes 2.0--6.0 exchange events through the
Kafka-compatible messaging layer; they should not be shown as direct
synchronous calls unless such a call is implemented.

\begin{longtable}{|p{0.65in}|p{1.4in}|p{2in}|p{2.3in}|}
\hline
\textbf{Process} & \textbf{Name} & \textbf{Inputs} & \textbf{Outputs} \\
\hline
1.0 & Authenticate & Login credentials or refresh token &
Access/refresh tokens and user role \\
\hline
2.0 & Ingest and extract documents & Uploaded file and request details &
Classified/extracted document, vendor match, review-queue items \\
\hline
3.0 & Approve and reserve & Purchase request and approval/rejection
signal & Decision, stock/licence reservations, escalations \\
\hline
4.0 & Manage contracts & Approved request and signed callback &
Generated contract, signature status, renewal reminders \\
\hline
5.0 & Score vendor risk & Vendor data and outcomes &
Risk score, band, contributing factors, drift flag \\
\hline
6.0 & Notify & Events from processes 2.0--5.0 &
Emails, digests, notification log \\
\hline
\end{longtable}

\subsubsection{Sequence Diagram}

The sequence diagram should show the purchase-request journey from the
requester and frontend/gateway through document intake, Kafka, the
approval API and worker, Temporal, the contract service, and the
notification service, ending with the signed-contract event and
request-status update.

\subsubsection{Activity Diagram: Document Processing}

The document-processing activity flow begins with upload and malware
scanning. The upload is rejected if the file is flagged or the scanner
is unavailable. A clean file proceeds through parsing, classification,
field extraction, vendor matching, duplicate checking, and confidence
scoring. Results meeting the threshold can be completed; lower-confidence
results are sent to a human review queue.

\textbf{Figure 4.2:} Activity diagram of document processing. Insert
the regenerated diagram with readable decision branches and labels.

\subsubsection{Class/Entity Relationship Diagram}

The principal entities identified in the report are User, Vendor,
PurchaseRequest, Document, ExtractedField, ApprovalRecord,
InventoryItem, LicenseUsage, Contract, RiskScore, Notification, and
RenewalReminder. The diagram should show the relationships described in
the project: a requester raises purchase requests; requests refer to
vendors and can have documents and approval records; vendors can have
documents, contracts, and risk scores; documents have extracted fields;
contracts have renewal reminders; and inventory items are associated
with licence-usage records.

The exact table definitions are maintained in the shared database
schema and service migrations.

\textbf{Figure 4.3:} Entity relationship diagram. Insert the regenerated
diagram and check that cardinalities and relationship labels are
readable and consistent with the actual database schema.

\subsection{Cost Analysis}

The project uses open-source software, so the report identifies no
additional software-licence cost. Development is based on existing
student laptops. The principal potential costs are team effort, external
API usage beyond free tiers, optional cloud hosting, and a production
email/domain service. The estimates below are indicative and should be
updated if actual quotations or usage figures become available.

\begin{longtable}{|p{1.5in}|p{3.5in}|p{1.4in}|}
\hline
\textbf{Item} & \textbf{Basis} & \textbf{Indicative Cost} \\
\hline
Software licences & Python, FastAPI, PostgreSQL, Redis, Redpanda/Kafka,
MinIO, Temporal, MLflow, Prometheus, Grafana, ClamAV, Docling, React,
Docker and other listed open-source tools & Rs. 0 \\
\hline
Development hardware & Existing student laptops & Rs. 0 additional \\
\hline
External APIs & Free tiers where available; usage depends on provider
limits & Depends on usage \\
\hline
Team effort & Part-time project work by team members & Not monetised \\
\hline
Optional cloud hosting & Pilot virtual machine with approximately
8 vCPU and 16 GB RAM, plus storage and backups &
Approx. Rs. 6,000--10,000 per month \\
\hline
Domain and email delivery & Production SMTP service in place of
development mail tooling & A few hundred rupees per month at pilot
volume \\
\hline
\end{longtable}

A meaningful break-even estimate requires measured monthly document
volume, average licence costs, and an estimate of time saved or duplicate
payments avoided. Those figures are not specified here, so no
return-on-investment claim is made.

\subsection{Hardware and Software Requirements}

The stack runs multiple containers concurrently. The document parser
and ClamAV can be resource-intensive; therefore, the recommended
configuration is expected to provide a smoother development and
demonstration experience.

\begin{longtable}{|p{1.5in}|p{2in}|p{2.5in}|}
\hline
\textbf{Component} & \textbf{Minimum} & \textbf{Recommended} \\
\hline
Processor & 4 cores & 8 cores \\
\hline
Memory & 8 GB RAM & 16 GB RAM \\
\hline
Storage & 25 GB free & 50 GB free for images, model files, and
database volumes \\
\hline
Network & Required for initial builds, image pulls, and ClamAV
definition updates & Broadband connection \\
\hline
GPU & Not required & Optional for heavier OCR or transformer workloads \\
\hline
\end{longtable}

\begin{longtable}{|p{1.7in}|p{4.3in}|}
\hline
\textbf{Category} & \textbf{Requirement} \\
\hline
Operating system & macOS, Ubuntu/Linux, or Windows 10/11 with WSL2 \\
\hline
Container runtime & Docker Engine or Docker Desktop with Docker Compose \\
\hline
Backend & Python with FastAPI, SQLAlchemy, Alembic, pytest,
scikit-learn, rapidfuzz, and Jinja2 \\
\hline
Frontend & React, TypeScript, Vite, and Tailwind CSS; the frontend is
built in a Node 20 container \\
\hline
Data and messaging & PostgreSQL, Redis, Redpanda (Kafka API), and MinIO \\
\hline
Workflow and ML & Temporal and MLflow \\
\hline
Document processing & Docling, pdfplumber, Tesseract OCR, and ClamAV \\
\hline
Monitoring and mail & Prometheus, Grafana, and Mailpit \\
\hline
Gateway & Nginx \\
\hline
Browser & A current version of Chrome, Edge, Firefox, or Safari \\
\hline
\end{longtable}
\section{System Architecture}

The Agentic AI-Based IT Procurement Intelligence Platform follows a
modular, event-driven architecture. The system is divided into
independent services, each responsible for a specific procurement
function. This separation improves maintainability, supports parallel
development, and allows background tasks to be processed independently
of user-facing requests.

The frontend communicates with backend services through a common API
gateway. Synchronous HTTP requests are used for operations that require
an immediate response, while inter-service communication primarily
uses a Kafka-compatible messaging layer. Temporal supports durable
workflows such as approvals, escalations, and contract-renewal
processes.

PostgreSQL stores application records, Redis supports caching and
locking, and MinIO stores uploaded documents. Background workers
process events and execute tasks that may take longer to complete.
This design allows the application to return responses to users
without waiting for every downstream operation to finish.

\subsection{UI/UX Design}

The frontend is implemented as a React and TypeScript single-page
application using Vite and Tailwind CSS. Nginx serves the frontend
and routes API requests to the appropriate backend services. Serving
the frontend and API through a common entry point simplifies access
to the application.

The interface is designed around the responsibilities of the four
implemented roles: requester, approver, finance, and administrator.

The following design principles guide the user interface:

\begin{enumerate}[label=\arabic*.]
    \item \textbf{Role-aware navigation:}
    The interface provides navigation appropriate to requesters,
    approvers, finance users, and administrators. Frontend route
    checks improve usability but do not replace backend
    authorization.

    \item \textbf{Clear asynchronous status:}
    The interface communicates the status of document classification,
    extraction, approvals, and other operations that may take time
    to complete.

    \item \textbf{Transparent system limitations:}
    The interface should distinguish successful operations from
    pending or failed operations. It should not indicate that an
    action succeeded when the backend is unavailable.

    \item \textbf{Human review:}
    Users can review extracted document fields and their confidence
    scores. Where correction is supported, users can correct the
    extracted information before proceeding.

    \item \textbf{Demonstration support:}
    The login interface supports demonstration through demo-role
    switching, allowing the different role-specific workflows to
    be presented during project evaluation.
\end{enumerate}

The principal screens and their purposes are summarised in
Table 4.3.

\begin{longtable}{|p{1.35in}|p{1.35in}|p{3in}|}
\hline
\textbf{Screen} & \textbf{User Role} & \textbf{Purpose} \\
\hline
Login & All roles & Sign in or select a demonstration role. \\
\hline
Requester Dashboard & Requester & View submitted purchase requests and their current status. \\
\hline
New Request Wizard & Requester & Enter request details, select a vendor, optionally upload a document, review extracted fields, and submit the request. \\
\hline
Approval Inbox & Approver, Finance & Review pending requests and approve or reject them according to assigned permissions. \\
\hline
Document Review & Approver, Finance, Admin & Review extracted document fields and confidence scores, and correct and save fields where supported. \\
\hline
Vendor Profile & Finance, Admin & View vendor risk information and the payment-change verification queue, along with supported vendor actions. \\
\hline
Contracts & Admin & Generate contracts, manage the signature workflow, and view contract-renewal timelines. \\
\hline
Inventory & Approver, Admin & View hardware stock and software-licence utilisation. \\
\hline
Notification Log & All, subject to permissions & View sent, queued, and failed notifications. \\
\hline
Global Search & All, subject to permissions & Search available purchase requests, vendors, contracts, and documents. \\
\hline
\end{longtable}

\textbf{Figure 4.4: Main-screen wireframes of the platform.}

Insert screenshots or labelled wireframes of the login screen,
new-request wizard, approval inbox, and vendor profile. Replace
this instruction with the actual figure before submitting the
final report.

\subsection{Block Diagram}

The block diagram represents the major components of the platform
and the relationships between them. It should show the frontend,
gateway, application services, background workers, messaging and
workflow components, data stores, and monitoring tools.

The main architectural layers are described in Table 4.4.

\begin{longtable}{|p{1.4in}|p{4.3in}|}
\hline
\textbf{Layer} & \textbf{Components and Responsibilities} \\
\hline
Clients & The browser-based React application and API clients access the platform through the API gateway on port 8080. \\
\hline
Gateway & Nginx serves the frontend and routes API paths to the appropriate backend services, providing a common entry point. \\
\hline
Application Services & The Document and Vendor Agent runs on port 8001; the Approval and Inventory Agent on port 8002; the Contract and Risk Agent on port 8003; the Notification Agent on port 8004; and the Authentication Service on port 8005. Authentication and role-based middleware protect application endpoints, except for specifically designated health and metrics endpoints. \\
\hline
Background Workers & Each domain service can use a background worker to consume events and perform asynchronous processing and workflow-related tasks. \\
\hline
Messaging and Orchestration & Redpanda provides a Kafka-compatible event-messaging interface. Temporal supports durable workflows for approvals, escalations, and contract renewals. \\
\hline
Data and Platform Services & PostgreSQL stores application records, Redis supports caching and locks, and MinIO stores documents. ClamAV scans files for malware, MLflow supports model tracking, and Mailpit supports development email testing. \\
\hline
Observability & Prometheus collects metrics and Grafana provides dashboards for monitoring request rates, latency, errors, and service availability. \\
\hline
\end{longtable}

\textbf{Figure 4.5: Block diagram of the proposed system architecture.}

Insert a clean block diagram based on the architecture described
above. Verify that the service names, ports, gateway routes,
workers, event bus, workflow engine, data stores, and monitoring
components correspond to the implemented system.

\subsection{Main Data Paths}

The principal data paths in the system are described below.

\begin{enumerate}[label=\arabic*.]
    \item \textbf{User access:}
    The browser accesses the frontend through Nginx. API requests
    are routed to the appropriate backend service.

    \item \textbf{Document processing:}
    The Document and Vendor Agent scans an uploaded file for
    malware. Accepted files are stored in MinIO, and an ingestion
    event is published. Background processing extracts and
    classifies document information and publishes the resulting
    classification and vendor-matching information.

    \item \textbf{Approval processing:}
    The Approval and Inventory Agent interacts with the workflow
    engine to initiate or progress approval workflows. Background
    workers execute workflow-related tasks, apply decisions, and
    publish relevant events.

    \item \textbf{Contract processing:}
    The Contract and Risk Agent processes approved requests,
    generates contracts, and handles signed-document callbacks.
    The callback is accepted only after the required HMAC
    signature validation and replay checks. A contract-signed
    event is published, allowing the Approval and Inventory
    Agent to update the purchase-request status.

    \item \textbf{Risk assessment and notifications:}
    The Contract and Risk Agent publishes relevant risk-score
    and renewal events. The Notification Agent consumes
    applicable events and sends email notifications through
    the configured SMTP service.

    \item \textbf{Data access and supporting services:}
    Application services use PostgreSQL for persistent records.
    Redis supports approval and contract locking, MinIO stores
    documents, and MLflow supports tracking of contract and
    risk-model experiments.
\end{enumerate}

These data paths illustrate how synchronous API operations,
event-driven communication, and durable workflows work together
to support procurement activities.

\section{Chapter 4 Review Checklist Before Submission}

Before submitting this chapter for final review, verify the following
items:

\begin{enumerate}[label=\arabic*.]
    \item Replace every instruction to insert a diagram, wireframe,
    or screenshot with the actual figure, or remove the placeholder
    if the figure is not included in the final scope.

    \item Ensure that figure numbers and table numbers are unique
    and sequential throughout the report. Check consistency with
    the numbering used in Chapter 3.

    \item Verify service names, ports, user roles, event names,
    software versions, and system behaviour against the actual
    implementation.

    \item Recheck the cost estimates and any statements about
    external services, free-tier limits, and third-party
    integrations. Clearly identify estimates that have not been
    confirmed through actual quotations or usage data.

    \item Confirm that the functional and non-functional
    requirements describe implemented features accurately and
    distinguish implemented functionality from planned or
    externally dependent features.

    \item Update the table of contents, list of figures, and
    list of tables after integrating the chapter into the
    complete report.
\end{enumerate}

\end{document}